\chapter{Telegram-бот}
\label{ch:telegram}

Бот решает две задачи: пускает на сайт без пароля и собирает лобби прямо в чате партии.

\section{Начало работы}
\label{sec:telegram:start}

Напишите боту \texttt{/start} в личные сообщения. Это обязательный шаг: пока вы не начали диалог, Telegram запрещает боту писать вам первым, и все ссылки для входа до вас просто не дойдут.

\section{Вход в браузере}
\label{sec:telegram:link}

Команда \texttt{/link} \textbf{в личке} присылает одноразовую ссылку для входа. Перешли по ней --- вы авторизованы, пароль не нужен.

\begin{itemize}
\item Ссылка одноразовая: повторный переход по той же не сработает, запросите новую.
\item В группе команда не работает --- ссылка входа приходит только в личные сообщения.
\item Открывайте ссылку \textbf{в обычном браузере}, а не во встроенном браузере Telegram: во встроенном часть возможностей работает хуже.
\end{itemize}

\section{Создание лобби из чата}
\label{sec:telegram:lobby}

Команда вызывается \textbf{в группе} с игроками:

\begin{quote}
\ttfamily
/создать\_лобби Вечерняя партия @player1 @player2
\end{quote}

Первое слово после команды и до первого упоминания --- название лобби, дальше перечисляются игроки через \texttt{@username}. Команда понимает несколько написаний: \texttt{/создать\_лобби}, \texttt{/create\_lobby}, \texttt{/createlobby}, \texttt{/lobby}.

Что произойдёт: бот создаст лобби, а каждому упомянутому игроку пришлёт в личку одноразовую ссылку, ведущую сразу в это лобби. Игроку не нужно ни искать сайт, ни вводить пароль --- один переход, и он в комнате ожидания.

\paragraph{Требования}
\begin{itemize}
\item У вызвавшего команду должны быть права мастера (см.\ раздел \ref{sec:getting-started:master-rights}).
\item Каждый упомянутый игрок должен хотя бы раз написать боту \texttt{/start}, иначе ссылка до него не дойдёт. Бот сообщит, кому не смог написать.
\item У игрока должен быть заполнен \texttt{@username} в Telegram --- по нему бот их и находит.
\end{itemize}

\section{Если бот молчит}
\label{sec:telegram:silent}

\begin{itemize}
\item Проверьте, что писали \texttt{/link} именно в личку, а не в группу.
\item Если бот раньше отвечал, а теперь нет --- напишите \texttt{/start} заново.
\item При создании лобби бот отдельно перечисляет тех, кому не смог отправить ссылку. Этим людям нужно начать диалог с ботом и попросить мастера повторить команду.
\end{itemize}
